\chapter{Shift Operations}
\label{ch:shift-operations}

\section{Overview}

The SIRC-1 CPU supports powerful shift operations that can be applied to operands in many instructions. Shifts can
multiply or divide values by powers of 2, extract bit fields, or perform bit rotations.

Shift operations are encoded within Short Immediate and Register format instructions, allowing a value to be shifted
before being used in an ALU operation—all within a single instruction.

Shift operations are only applied to source operands. A shift cannot be applied to the result of an operation before
register writeback.

\textbf{Note:} By default, when shifts are used with ALU instructions (apart from the \mnemonic{SHFT} meta-instruction), the status register flags are updated based on the ALU operation result, not the shift. This behavior can be overridden using the status override syntax (see Section~\ref{sec:status-override}).

\section{Shift Types}

\subsection{Shift Type Encoding}

\begin{table}[H]
	\centering
	\begin{tabular}{clcp{7cm}}
		\toprule
		\textbf{Code} & \textbf{Mnemonic} & \textbf{Type}          & \textbf{Description}                               \\
		\midrule
		000           & NUL               & --                     & No shift performed                                 \\
		001           & LSL               & Logical Shift Left     & Shift left, fill with zeros                        \\
		010           & LSR               & Logical Shift Right    & Shift right, fill with zeros                       \\
		011           & ASL               & Arithmetic Shift Left  & Same data result as LSL; sets V on signed overflow \\
		100           & ASR               & Arithmetic Shift Right & Shift right, preserving sign bit                   \\
		101           & RTL               & Rotate Left            & Rotate left within the operand                     \\
		110           & RTR               & Rotate Right           & Rotate right within the operand                    \\
		111           & --                & Reserved               & Reserved for future use                            \\
		\bottomrule
	\end{tabular}
	\caption{Shift Type Encoding}
	\label{tab:shift-encoding}
\end{table}

\mnemonic{NUL} is the no-shift form: it leaves the operand unchanged but still routes it through the shifter, so \texttt{SHFT rD, NUL \#0} is a legal way to drive the status flags from the unshifted value of \texttt{rD}.

\section{Shift Syntax}

A shift suffix may follow the ALU short-immediate and ALU register forms, the \mnemonic{SHFT} meta-instruction, and the
register-displacement \mnemonic{LOAD}/\mnemonic{STOR} forms. It is rejected on \mnemonic{LOAD rD, rS}, on
\mnemonic{LDEA}/\mnemonic{LDEL}, and on coprocessor instructions.

For ALU instructions, the shift is applied in the "Decode and Register Fetch" phase, before the ALU operation. In
register-format instructions the shifted operand is R2 (first source); in short-immediate instructions it is the
register named by the Reg field (first source, since the short-immediate format has no separate source register). In
register-displacement \mnemonic{LOAD} forms the shift is applied to the offset register used to compute the effective
address, never to the value loaded from memory; in register-displacement \mnemonic{STOR} forms it is applied to the
source register before its value is written to memory.

If using the three-operand form of an instruction, the shifted operand is the second operand. If using the two-operand
form, the second operand is inferred to be the same as the destination.

See this table for some examples of which operand is shifted:

\begin{table}[H]
	\centering
	\begin{tabular}{lll}
		\toprule
		\textbf{Instruction}     & \textbf{How the CPU Interprets It} & \textbf{First Source Operand} \\
		\midrule
		ADDR r1, r2, r3, LSL \#1 & ADDR r1, r2, r3, LSL \#1           & r2                            \\
		ADDR r1, r2, LSL \#1     & ADDR r1, r1, r2, LSL \#1           & r1                            \\
		ADDI r1, \#2, LSL \#1    & ADDI r1, r1, \#2, LSL \#1          & r1                            \\
		\bottomrule
	\end{tabular}
	\caption{Examples of first source operands}
	\label{tab:source-operand-examples}
\end{table}

\section{Shift Type Details}

\subsection{Logical Shift Left (LSL)}

\textbf{Operation:} Shift bits to the left, filling vacated bits with zeros.

\begin{figure}[H]
	\centering
	\begin{bytefield}[bitwidth=1.2em]{17}
		\bitheader{0-15} \\
		\bitbox{1}{$\leftarrow$}
		\bitbox{15}{\small Shift Left}
		\bitbox{1}{0}
	\end{bytefield}
	\caption{Logical Shift Left}
\end{figure}

\begin{itemize}
	\item The last bit shifted out of bit 15 (bit 16-n for a count of n) is placed in the carry flag
	\item Bit 0 receives 0
	\item Equivalent to multiplying by $2^n$ where $n$ is the shift count
\end{itemize}

\textbf{Example 9-1:}
\begin{lstlisting}
; Multiply r1 by 4
ADDI r1, #0, LSL #2  ; r1 = (r1 << 2) = r1 * 4

; Set bit 10 (16-bit immediate form)
ORRI r3, #0x0400         ; r3 |= 0x0400
\end{lstlisting}

\subsection{Logical Shift Right (LSR)}

\textbf{Operation:} Shift bits to the right, filling vacated bits with zeros.

\begin{figure}[H]
	\centering
	\begin{bytefield}[bitwidth=1.2em]{17}
		\bitheader{0-15} \\
		\bitbox{1}{0}
		\bitbox{15}{\small Shift Right}
		\bitbox{1}{$\rightarrow$}
	\end{bytefield}
	\caption{Logical Shift Right}
\end{figure}

\begin{itemize}
	\item The last bit shifted out of bit 0 (bit n-1 for a count of n) is placed in the carry flag
	\item Bit 15 receives 0
	\item Equivalent to unsigned division by $2^n$
	\item Use for unsigned values only
\end{itemize}

\textbf{Example 9-2:}
\begin{lstlisting}
; Divide unsigned value by 8
ADDI r1, #0, LSR #3  ; r1 = (r1 >> 3) = r1 / 8 (unsigned)

; Extract high byte
ADDI r3, #0, LSR #8  ; r3 = (r3 >> 8)
\end{lstlisting}

\subsection{Arithmetic Shift Left (ASL)}

\textbf{Operation:} Produces the same data result as LSL, but sets the overflow flag if the sign bit changes.

\begin{itemize}
	\item Equivalent to signed multiplication by $2^n$
	\item Can cause signed overflow if sign bit changes
	\item Overflow flag is set if sign bit changes during shift
\end{itemize}

\textbf{Example 9-3:}
\begin{lstlisting}
; Multiply signed value by 2
ADDI r1, #0, ASL #1  ; r1 = (r1 << 1) = r1 * 2 (signed)
\end{lstlisting}

\subsection{Arithmetic Shift Right (ASR)}

\textbf{Operation:} Shift bits to the right, preserving the sign bit (bit 15).

\begin{figure}[H]
	\centering
	\begin{bytefield}[bitwidth=1.2em]{17}
		\bitheader{0-15} \\
		\bitbox{1}{$\curvearrowleft$}
		\bitbox{15}{\small Shift Right}
		\bitbox{1}{$\rightarrow$}
	\end{bytefield}
	\caption{Arithmetic Shift Right}
\end{figure}

\begin{itemize}
	\item The last bit shifted out of bit 0 (bit n-1 for a count of n) is placed in the carry flag
	\item Bit 15 (sign bit) is replicated into all bits vacated by the shift
	\item Equivalent to signed division by $2^n$ (rounds toward negative infinity)
	\item Use for signed values only
\end{itemize}

\textbf{Example 9-4:}
\begin{lstlisting}
; Divide signed value by 4
ADDI r1, #0, ASR #2  ; r1 = (r1 >> 2) = r1 / 4 (signed)

; Average two signed values
ADDR r3, r4, r5
ADDI r3, #0, ASR #1  ; r3 = (r3 >> 1) = (r4 + r5) / 2
\end{lstlisting}

\subsection{Rotate Left (RTL)}

\textbf{Operation:} Rotate bits to the left within the 16-bit operand. The bit that wraps from bit 15 to bit 0 is also
copied into the carry flag. The incoming carry flag is not used as an input.

\begin{figure}[H]
	\centering
	\begin{bytefield}[bitwidth=1.2em]{17}
		\bitheader{0-15} \\
		\bitbox{1}{$\curvearrowleft$}
		\bitbox{15}{\small Rotate Left}
		\bitbox{1}{$\curvearrowleft$}
	\end{bytefield}
	\caption{Rotate Left}
\end{figure}

\begin{itemize}
	\item Bit 15 wraps into bit 0 and is copied into the carry flag
	\item The previous carry flag value is ignored
	\item No data is lost
	\item Useful for bit manipulation, circular buffers, and cyclic bit patterns
\end{itemize}

\textbf{Example 9-5:}
\begin{lstlisting}
; Rotate r1 left by 1
SHFT r1, RTL #1

; Circular rotate of one word
SHFT r3, RTL #4
\end{lstlisting}

\subsection{Rotate Right (RTR)}

\textbf{Operation:} Rotate bits to the right within the 16-bit operand. The bit that wraps from bit 0 to bit 15 is also
copied into the carry flag. The incoming carry flag is not used as an input.

\begin{figure}[H]
	\centering
	\begin{bytefield}[bitwidth=1.2em]{17}
		\bitheader{0-15} \\
		\bitbox{1}{$\curvearrowright$}
		\bitbox{15}{\small Rotate Right}
		\bitbox{1}{$\curvearrowright$}
	\end{bytefield}
	\caption{Rotate Right}
\end{figure}

\begin{itemize}
	\item Bit 0 wraps into bit 15 and is copied into the carry flag
	\item The previous carry flag value is ignored
	\item No data is lost
\end{itemize}

\textbf{Example 9-6:}
\begin{lstlisting}
; Rotate r1 right by 1
SHFT r1, RTR #1

; Check if bit 0 is set
SHFT r2, RTR #1      ; Bit 0 -> carry
; Now branch on carry to test the old bit 0
\end{lstlisting}

\section{Shift Operands}

The Shift Operand (SO) bit determines what value is shifted:

\begin{table}[H]
	\centering
	\begin{tabular}{clp{8cm}}
		\toprule
		\textbf{SO Bit} & \textbf{Mode}           & \textbf{Description}                                                                                                                       \\
		\midrule
		0               & Shift by Immediate      & The first source operand (Register format: R2; Short Immediate format: the Reg field) is shifted by a literal amount                       \\
		1               & Shift Count in Register & The first source operand (Register format: R2; Short Immediate format: the Reg field) is shifted by the amount in the shift count register \\
		\bottomrule
	\end{tabular}
	\caption{Shift Operand Modes}
	\label{tab:shift-operand}
\end{table}

\subsection{Mode 0: Shift by Immediate}

The first source operand (R2) is shifted by a literal amount before the ALU operation.

\begin{lstlisting}
; Multiply r1 by 4 and add 1
ADDI r1, #1, LSL #2  ; r1 = (r1 << 2) + 1
                          ; SO=0: r1 is shifted by literal 2

\end{lstlisting}

\subsection{Mode 1: Shift by Register Count}

The shift count is taken from the shift amount field (interpreted as a register), allowing variable-length shifts. R2
(first source) is still what gets shifted.

\begin{lstlisting}
; Variable shift
; Shift amount register contains shift count
ADDR r1, r2, r3, LSL r4  ; r1 = (r2 << r4) + r3
                          ; SO=1: shift r2 by value in r4

; Barrel shifter emulation
ADDR r4, r5, r6, ASR r7  ; r4 = (r5 >> r7) + r6 (arithmetic)
\end{lstlisting}

\section{Shift Count}

The shift count is 0--15 when specified as an immediate (SO = 0); the field is 4 bits wide, and both the assembler and
the encoder enforce this range, so a count greater than 15 can never occur through the immediate path. When the shift
count is taken from a register (SO = 1), the full 16-bit register value is available as a count, and a count greater
than 15 is well defined:

\begin{itemize}
	\item Shift count of 0 means no shift is performed
	\item Maximum immediate shift count is 15
	\item For \mnemonic{LSL}, \mnemonic{LSR}, and \mnemonic{ASL}, a register-supplied count of 16 or more produces the
	      same result as a count of exactly 16: every original bit is shifted out, giving \texttt{0x0000}. The carry
	      flag reflects the last bit shifted out at count 16, as for any other count.
	\item For \mnemonic{ASR}, a register-supplied count of 16 or more sign-extends fully: \texttt{0xFFFF} if the
	      original value was negative, \texttt{0x0000} if it was positive or zero. This is the same rule as any other
	      count, just carried to its limit -- every vacated bit is the sign bit.
	\item \mnemonic{RTL} and \mnemonic{RTR} have no vacated bits to define, since a rotate never discards data; a
	      register-supplied count is used modulo 16, so a count of 16 behaves as a count of 0, a count of 17 behaves
	      as a count of 1, and so on.
\end{itemize}

\section{Status Flag Effects}

\subsection{Default Behavior}

When shift operations are used with ALU instructions, the status flags are \textbf{by default} updated based on the
\textbf{ALU operation result}, not the shift operation. For example, in \texttt{CMPI r2, \#1, LSL \#2}, the flags
reflect the comparison result, not the shift.

To update flags based on the shift operation instead, use the \texttt{[S]} status override modifier (see
Section~\ref{sec:status-override}), or use the \mnemonic{SHFT} meta-instruction documented in
Chapter~\ref{ch:alu-instructions}.

\subsection{Carry Flag}

For all shifts and rotates:
\begin{itemize}
	\item The last bit shifted out is placed in the carry flag
	\item For shifts of more than 1 position, only the final bit shifted out matters
	\item A shift count of 0 leaves the operand unchanged and clears the carry flag. When the status register is updated from the
	      shifter, all four condition flags are written, so no flag survives a shift-sourced update; to preserve flags across a
	      shift, use the \texttt{[N]} status override
\end{itemize}

\subsection{Zero Flag}

Set if the result after shifting is zero.

\subsection{Negative Flag}

Set if bit 15 of the result is 1 (result is negative when interpreted as signed).

\subsection{Overflow Flag}

\begin{itemize}
	\item For \textbf{ASL}: Set if the sign bit changes during the shift (signed overflow)
	\item For other shifts: Cleared to 0
\end{itemize}

\section{Common Use Cases}

\subsection{Multiplication by Constants}

\textbf{Example 9-7:}
\begin{lstlisting}
; Multiply by 2
ADDI r1, #0, LSL #1   ; r1 = (r1 << 1) = r1 * 2

; Multiply by 3
ADDR r1, r2, r2, LSL #1   ; r1 = (r2 << 1) + r2 = r2 * 3

; Multiply by 5
ADDR r1, r2, r2, LSL #2   ; r1 = (r2 << 2) + r2 = r2 * 5
\end{lstlisting}

\subsection{Division by Powers of 2}

\textbf{Example 9-8:}
\begin{lstlisting}
; Unsigned divide by 4
LOAD r1, r2        ; r1 = r2
ADDI r1, #0, LSR #2   ; r1 = (r1 >> 2) = r1 / 4 (unsigned)

; Signed divide by 8
LOAD r1, r2        ; r1 = r2
ADDI r1, #0, ASR #3   ; r1 = (r1 >> 3) = r1 / 8 (signed)
\end{lstlisting}

\subsection{Bit Field Extraction}

\textbf{Example 9-9:}
\begin{lstlisting}
; Extract bits 8-11 from r1
LOAD r2, r1        ; r2 = r1
ADDI r2, #0, LSR #8   ; Shift down
ANDI r2, #0x0F            ; Mask to 4 bits

; Extract and sign-extend bit 7
LOAD r2, r1        ; r2 = r1
ADDI r2, #0, LSL #8   ; Shift bit 7 to bit 15
ADDI r2, #0, ASR #8   ; Arithmetic shift to sign-extend
\end{lstlisting}

\subsection{Bit Manipulation}

\textbf{Example 9-10:}
\begin{lstlisting}
; Set bit N (where N is in r2)
LOAD r3, #1               ; r3 = 1
ADDI r3, #0, LSL r2       ; r3 = (r3 << r2) = 1 << N
ORRR r1, r3               ; r1 |= (1 << N)

; Clear bit N
LOAD r3, #1
ADDI r3, #0, LSL r2       ; r3 = (r3 << r2) = 1 << N
LOAD r4, #0xFFFF
XORR r3, r4, r3           ; r3 = 0xFFFF ^ (1 << N) = ~(1 << N)
ANDR r1, r1, r3           ; r1 = r1 & ~(1 << N)

; Test bit N
LOAD r3, #1
ADDI r3, #0, LSL r2       ; r3 = (r3 << r2) = 1 << N
TSAR r1, r3               ; Test r1 & (1 << N)
BRAN|!= @bit_was_set      ; Branch if bit set
\end{lstlisting}

\section{Multi-Word Shifts}

For operations on values larger than 16 bits, shifts and rotates can be combined:

\subsection{32-bit Left Shift}

\textbf{Example 9-11:}
\begin{lstlisting}
; Shift 32-bit value in r1:r2 left by 1
; r1 = high word, r2 = low word
SHFT r2, LSL #1       ; Shift low, MSB -> carry
ADCI r1, #0, LSL #1   ; Shift high and add carry into bit 0
\end{lstlisting}

\subsection{32-bit Right Shift (Unsigned)}

\textbf{Example 9-12:}
\begin{lstlisting}
; Shift 32-bit value in r1:r2 right by 1 (unsigned)
SHFT r1, LSR #1       ; Shift high, LSB -> carry
ORRI[N] r2, #0, LSR #1    ; Shift low while preserving carry from high word
ORRI|CS r2, #0x8000   ; If old high bit 0 was set, insert it into low bit 15
\end{lstlisting}

\subsection{32-bit Right Shift (Signed)}

\textbf{Example 9-13:}
\begin{lstlisting}
; Shift 32-bit value in r1:r2 right by 1 (signed)
SHFT r1, ASR #1       ; Arithmetic shift high, LSB -> carry
ORRI[N] r2, #0, LSR #1    ; Shift low while preserving carry from high word
ORRI|CS r2, #0x8000   ; If old high bit 0 was set, insert it into low bit 15
\end{lstlisting}

\section{Performance Considerations}

\begin{itemize}
	\item All shift operations complete in the same number of cycles regardless of shift count
	\item Shifts can be performed "for free" as part of another instruction (no extra cycles)
	\item Using shifts for multiplication/division by powers of 2 is much faster than using a multiply/divide instruction (if one
	      existed)
	\item Barrel shifter implementation in hardware allows any shift count in constant time
\end{itemize}

\section{Limitations}

\begin{itemize}
	\item The immediate shift count is limited to 0--15 (4-bit field)
	\item A register-supplied shift count greater than 15 is architecturally undefined
	\item For a shift greater than 15 positions, use multiple shift instructions with immediate counts
	\item Rotates do not consume the carry flag; \mnemonic{RTL} and \mnemonic{RTR} rotate only within the 16-bit operand
	\item The reserved shift type (111) has architecturally undefined behavior
\end{itemize}

\section{Status Flag Update Control}
\label{sec:status-override}

By default, ALU instructions with shift operations update the status register flags based on the \textbf{ALU operation
	result}, not the shift operation. However, the SIRC-1 provides manual control over which operation updates the flags.

\subsection{Status Override Syntax}

A status override modifier can be added immediately after an instruction mnemonic using square brackets:

\begin{table}[H]
	\centering
	\begin{tabular}{clp{8cm}}
		\toprule
		\textbf{Modifier} & \textbf{Meaning} & \textbf{Description}                    \\
		\midrule
		{[}A{]}           & ALU              & Update flags from ALU result (default)  \\
		{[}S{]}           & Shift            & Update flags from shift result, not ALU \\
		{[}N{]}           & None             & Do not update status flags at all       \\
		\bottomrule
	\end{tabular}
	\caption{Status Flag Override Modifiers}
	\label{tab:status-override}
\end{table}

\subsection{Examples}

\textbf{Example 9-14:}
\textbf{Default Behavior (ALU flags):}
\begin{lstlisting}
; Without override, flags reflect the ALU addition result
ADDI r2, #1, LSL #2           ; r2 = (r2 << 2) + 1
                              ; Flags: based on addition result
\end{lstlisting}

\textbf{Example 9-15:}
\textbf{Shift Flag Override:}
\begin{lstlisting}
; Update flags based on shift result instead of ALU
ADDI[S] r2, #1, LSL #2        ; r2 = (r2 << 2) + 1
                              ; Flags: based on (r2 << 2), not the +1

; Check if shifted value is zero
CMPI[S] r3, #0, LSR #4        ; Compare (r3 >> 4) with 0
                              ; Flags: based on shift, not comparison
\end{lstlisting}

\textbf{Example 9-16:}
\textbf{No Flag Update:}
\begin{lstlisting}
; Perform operation without affecting status flags
ADDI[N] r1, #5, LSL #1        ; r1 = (r1 << 1) + 5
                              ; Flags: unchanged

; Useful when preserving flags from previous operation
CMPR r4, r5                   ; Set flags based on comparison
ADDI[N] r6, #1                ; Increment r6 without changing flags
BRAN|== @equal_branch         ; Branch using flags from CMPR
\end{lstlisting}

\subsection{Use Cases}

\subsubsection{Testing Shift Results}

To determine properties of a shifted value without performing a subsequent comparison:

\begin{lstlisting}
; Check if r1 << 3 would be negative
ADDI[S] r1, #0, LSL #3        ; r1 = r1 << 3, flags from shift
BRAN|NS @overflow_detected    ; Branch if sign bit set
\end{lstlisting}

\subsubsection{Preserving Flags}

When performing operations that should not affect condition codes:

\begin{lstlisting}
; Test condition
CMPR r1, r2                   ; Compare r1 and r2

; Do some work without affecting flags
ADDI[N] r3, #1                ; Increment counter (flags unchanged)
LOAD r4, (#0, a)+             ; Load next value (LOAD never updates flags)

; Branch using original comparison flags
BRAN|>= @r1_was_greater
\end{lstlisting}

\subsubsection{SHFT Meta-Instruction}

The \mnemonic{SHFT} meta-instruction is the ALU-family shorthand for this common case:

\begin{lstlisting}
SHFT r1, LSL #3               ; ORRI[S] r1, #0, LSL #3
\end{lstlisting}

See Chapter~\ref{ch:alu-instructions} for the full \mnemonic{SHFT} instruction entry.

\subsection{Assembler Syntax Notes}

\begin{itemize}
	\item The status override modifier must appear immediately after the instruction mnemonic with no spaces: \texttt{ADDI[S]}
	\item The modifier is optional; omitting it defaults to {[}A{]} (ALU flag updates)
	\item The override applies to the entire instruction, including any conditional execution
	\item Explicit {[}A{]} is rarely needed but can improve code readability
\end{itemize}
